% ID: SEL-D0041-OB-E02
% ESTADO: cotejada. Segunda pasada visual y borrador compilado revisados.
% Fuente: Selectividad/matematicas_aplicadas/Ua6MATEMÁTICAS APLICADAS examen.pdf
% SHA-256: b7541cc45d3e96297c29037fe17a562b7daa9cd0ee2a95c2b5e9454ad9b90f14
% Página: 2; opción B; ejercicio original 2.

Se considera la función $f$, definida a trozos por la expresión
\[
f(x)=\begin{cases}
-x^2+x+6 & \text{si }x\leq2\\[4pt]
x+2 & \text{si }x>2.
\end{cases}
\]
\begin{enumerate}[label=\alph*)]
\item \textbf{(0.5 puntos)} Estudie la continuidad de la función.
\item \textbf{(0.5 puntos)} Analice la derivabilidad de la función.
\item \textbf{(1.5 puntos)} Represéntela gráficamente, determinando los extremos, los intervalos de crecimiento y decrecimiento y los puntos de corte con los ejes.
\end{enumerate}
